\section{Erd\H{o}s Problem \#313}

\subsection*{1) ROUND-2 OBJECTIVE}
Path pursued: \textbf{(C) obstruction/fail-safe mode}.  In Round~1 the problem was reduced to the infinitude of primary pseudoperfect numbers.  The remaining gap is the absence of sharp structural restrictions on a hypothetical infinite family.  In this round I do not re-prove the Round~1 reduction; instead I prove new congruence constraints that every solution must satisfy.  These constraints are genuinely stronger than the Round~1 output and isolate a precise Euclid-type bottleneck that any recursive construction must pass through.

\subsection*{2) Round-1 FOUNDATION USED}
I use the following Round~1 facts without re-proving them.
\begin{enumerate}
\item Any solution of
\[
\frac1{p_1}+\cdots+\frac1{p_k}=1-\frac1m
\]
forces
\[
 m=\prod_{i=1}^k p_i,
\]
so the problem is equivalent to the infinitude of squarefree integers $m>1$ satisfying
\[
\frac1m+\sum_{p\mid m}\frac1p=1.
\]
\item The numbers $m$ satisfying this equation are the primary pseudoperfect numbers.
\item If $m$ is primary pseudoperfect and $m+1$ is prime, then $m(m+1)$ is primary pseudoperfect.
\end{enumerate}

\subsection*{3) NEW INSIGHT / TOOL (ROUND-2)}
The new tool is a local congruence lemma extracted from the integer form of the equation.

\paragraph{Local congruence lemma.}
If $m$ is primary pseudoperfect, then for every prime divisor $p$ of $m$,
\[
\frac{m}{p}\equiv -1 \pmod p.
\]
This has two immediate consequences that were absent from Round~1.
\begin{enumerate}
\item If $P$ is the largest prime divisor of $m$, then
\[
P\mid \left(\frac{m}{P}+1\right)=\left(\prod_{\substack{q\mid m \\ q<P}} q\right)+1.
\]
So the largest prime factor of any solution must come from a divisor of a Euclid number.
\item If $6\mid m$, then
\[
m\equiv 6 \pmod{36}.
\]
Thus every $6$-divisible solution lies in a single residue class modulo $36$.
\end{enumerate}
These are real structural obstructions: they sharply restrict any inductive construction, and they explain why the naive ``keep appending an arbitrary prime'' strategy cannot work.

\subsection*{4) ATTACK PLAN (ROUND-2)}
The exact Round~1 gap was this: after reducing to primary pseudoperfect numbers, there was still no fine-grained arithmetic information on the prime factors of a solution.  The plan is therefore:
\begin{enumerate}
\item rewrite the defining equation as an integer identity,
\item reduce it modulo each prime divisor,
\item extract the strongest immediate corollaries, especially for the largest prime divisor and for the residue class modulo $36$.
\end{enumerate}
This overcomes one of the main obstacles listed in Round~1: it gives a rigorous obstruction to broad classes of naive recursive constructions.

\subsection*{5) WORK (ROUND-2)}
By the Round~1 reduction, any solution is a squarefree integer $m>1$ satisfying
\[
\frac1m+\sum_{p\mid m}\frac1p=1.
\]
Multiplying by $m$ gives the integer identity
\[
1+\sum_{p\mid m}\frac{m}{p}=m.
\]

\paragraph{Lemma 5.1.}
For every prime divisor $p$ of $m$,
\[
\frac{m}{p}\equiv -1 \pmod p.
\]

\paragraph{Proof.}
Reduce
\[
1+\sum_{q\mid m}\frac{m}{q}=m
\]
modulo $p$.  Because $m$ is squarefree, every term $m/q$ with $q\neq p$ is divisible by $p$, while $m/p$ is not.  Also $m\equiv 0\pmod p$.  Hence
\[
1+\frac{m}{p}\equiv 0 \pmod p,
\]
which is exactly the claim.  \hfill $\square$

\paragraph{Corollary 5.2 (largest-prime Euclid divisor law).}
Let $P$ be the largest prime divisor of $m$.  Then
\[
P\mid \left(\frac{m}{P}+1\right)=\left(\prod_{\substack{q\mid m \\ q<P}} q\right)+1.
\]

\paragraph{Proof.}
Apply Lemma~5.1 with $p=P$.  Since $m$ is squarefree, $m/P$ is exactly the product of the smaller prime divisors of $m$.  \hfill $\square$

\paragraph{Corollary 5.3 (mod $36$ law under $6\mid m$).}
If $6\mid m$, then
\[
m\equiv 6 \pmod{36}.
\]

\paragraph{Proof.}
Write $m=6N$.
Applying Lemma~5.1 with $p=2$ gives
\[
\frac{m}{2}=3N\equiv -1 \pmod 2,
\]
so $N$ is odd.
Applying Lemma~5.1 with $p=3$ gives
\[
\frac{m}{3}=2N\equiv -1 \pmod 3.
\]
Since $-1\equiv 2\pmod 3$, this becomes $2N\equiv 2\pmod 3$, hence $N\equiv 1\pmod 3$.
Therefore $N$ is odd and congruent to $1$ modulo $3$, so
\[
N\equiv 1 \pmod 6.
\]
Multiplying by $6$ yields
\[
m=6N\equiv 6 \pmod{36}.
\]
\hfill $\square$

\paragraph{Interpretation.}
Corollary~5.2 shows that whenever one builds a solution by adjoining a new largest prime factor, that prime must divide a Euclid number associated to the already chosen smaller prime factors.  Thus the chain lemma from Round~1 is not merely a convenient construction: it reflects a genuine arithmetic bottleneck built into every solution.

\subsection*{6) ADVERSARIAL VERIFICATION}
I checked the new argument against the most likely failure modes.
\begin{enumerate}
\item \emph{Did I smuggle in the unproved claim $6\mid m$?}  No.  Corollary~5.3 is explicitly conditional: it says \emph{if} $6\mid m$, \emph{then} $m\equiv 6\pmod{36}$.  This avoids conflict with the still-open possibility of odd primary pseudoperfect numbers.
\item \emph{Is the largest-prime corollary using squarefreeness correctly?}  Yes.  Squarefreeness is a Round~1 foundation result, so $m/P$ is exactly the product of the other prime divisors.
\item \emph{Do the corollaries agree with the known examples?}  Yes:
\[
7\mid (2\cdot 3)+1,
\qquad
43\mid (2\cdot 3\cdot 7)+1,
\qquad
31\mid (2\cdot 3\cdot 11\cdot 23)+1=1519,
\]
and all known primary pseudoperfect numbers larger than $2$ are congruent to $6$ modulo $36$.
\end{enumerate}
No hidden quantifier issue remains: the lemma applies to \emph{every} prime divisor $p\mid m$.

\subsection*{7) FINAL}
\textbf{UNRESOLVED (BUT STRICTLY ADVANCED).}

Round~2 proves new necessary conditions that were not present in Round~1:
\begin{enumerate}
\item the local congruence
\[
\frac{m}{p}\equiv -1\pmod p \qquad (p\mid m),
\]
\item the largest-prime Euclid-divisor law
\[
P\mid \left(\prod_{q<P,\ q\mid m} q +1\right),
\]
\item the residue restriction
\[
6\mid m \Longrightarrow m\equiv 6\pmod{36}.
\]
\end{enumerate}
These sharpen the structural picture of primary pseudoperfect numbers and rule out broad classes of naive constructions.  The infinitude question itself remains open.

\subsection*{8) COMPLETION ESTIMATE}
COMPLETION: 35\%

\subsection*{9) REFERENCES}
\begin{enumerate}
\item T. F. Bloom, \emph{Erd\H{o}s Problem \#313}, problem page, accessed 2026-03-07.
\item OEIS A054377, \emph{Primary pseudoperfect numbers}, accessed 2026-03-07.
\item J. Sondow and K. MacMillan, \emph{Primary Pseudoperfect Numbers, Arithmetic Progressions, and the Erd\H{o}s--Moser Equation}, American Mathematical Monthly 124 (2017).
\end{enumerate}

\bigskip
\hrule
\bigskip

